\begin{exbox}[三边带截断的归一化与第一谐波]
沿用前面的 $b=|\beta|$，并记 $s=\omega_{\rm r}t=s_{\rm d}$。
三边带近似则保留 $c_0=J_0(2b)$、$c_{\pm1}=\pm J_1(2b)
\ee^{\pm\ii\varphi_\beta}$。其保留概率为
$S_3=J_0(2b)^2+2J_1(2b)^2$，遗漏概率为 $1-S_3$。
若条件于这三个通道，并令 $\theta=\kappa_z\zeta+\varphi_\beta$，周期密度是
\begin{equation}
\mathfrak n_3(\theta,s)=\frac{1}{S_3}
\left[J_0^2+4J_1^2\sin^2\theta
+4J_0J_1\sin s\sin\theta\right].
\label{eq:feqo-three-sideband-density}
\end{equation}
其第一 Fourier 系数的模长为 $2|J_0J_1\sin s|/S_3$，
正弦调制幅度是这个数的两倍。
三边带模型的极大在 $s=\pi/2$，但当遗漏概率不小，它不能替代完整 Bessel 梳的优化。
\end{exbox}

\begin{exbox}[相对论漂移长度与有限能散的数值核对]
取动能 $K_0=\SI{200}{\kilo\electronvolt}$、波长 $\lambda=\SI{800}{\nano\metre}$。
有 $\gamma_0=1.39139$、$v_0=2.08450\times10^8\,\si{\metre\per\second}$，
$\omega=2.35456\times10^{15}\,\si{\per\second}$，因此
\begin{equation}
L_{\rm T}=\frac{4\pi\gamma_0^3m_{\rm e}v_0^3}{\hbar\omega^2}
=\SI{0.47770}{\metre}.
\label{eq:feqo-talbot-numerical}
\end{equation}
半复现长度 $L_{\rm T}/2$ 是把密度平移半周期的长度；原位复现需要上式的完整长度。
对 $b=1.20$，第一谐波的最早理想极大在 $L_1^*=\SI{0.02993}{\metre}$。

若源的中心能量有标准差 $\sigma_E=\SI{1}{\electronvolt}$ 的独立高斯涨落，
$\partial_Ev=1/(\gamma_0^3m_{\rm e}v_0)$ 给
\[
\sigma_t(L)=\frac{L\sigma_E}{\gamma_0^3m_{\rm e}v_0^3},\qquad
D_q(L)=\exp[-q^2\omega^2\sigma_t(L)^2/2].
\]
定义 $L_E$ 为第一谐波衰减到 $\ee^{-1/2}$ 的长度，则
\begin{equation}
L_E=\frac{\gamma_0^3m_{\rm e}v_0^3}{\omega\sigma_E}
=\SI{0.05891}{\metre}.
\label{eq:feqo-LE-def}
\end{equation}
在上述理想极大处 $D_1=0.879$，而 $D_2=0.597$；高次谐波先被洗掉。
因此最佳实际距离应最大化
$|J_1(4b\sin(2\pi L/L_{\rm T}))|\exp[-L^2/(2L_E^2)]$，
它通常比无噪声最优距离更短。

若改取 $K_0=\SI{100}{\kilo\electronvolt}$、$\lambda=\SI{2}{\micro\metre}$、$b=1$，
同一计算给 $L_{\rm T}=\SI{0.92869}{\metre}$、$L_1^*=\SI{0.07070}{\metre}$。
光学周期为 $\SI{6.6713}{\femto\second}$，但一个谐波极大本身不能定义脉冲全宽；
全宽必须从完整周期密度求峰值和半高交点，不能从 $b$ 单独指定一个通用常数。
明确一种可复查的定义：先选定一个峰 $\tau_*$ 和背景值 $n_{\rm bg}$，
在该峰相邻的两个单调区间解
\begin{equation}
 n(\tau_\pm)=n_{\rm bg}+\frac{n(\tau_*)-n_{\rm bg}}2,\qquad
 W_{1/2}=\tau_+-\tau_-.
 \label{eq:feqo-pulse-width-definition}
\end{equation}
绝对半高宽取 $n_{\rm bg}=0$；相对背景的半对比宽取实际选定背景。
弱聚束若整个周期密度都高于绝对半峰值，则绝对半高交点不存在。
多峰、周期包络或不同背景约定也会改变结果；完整密度以及这个读出定义应一起报告。
\end{exbox}
